% -----------------------------------------------------------------------------
% This file is automatically generated.
% Do not edit manually.
% Source: notebooks/support/regression-analysis/support.Rmd
% -----------------------------------------------------------------------------

\par\addvspace{\topsep}
\begingroup
\fvset{listparameters={%
  \setlength{\topsep}{0pt}%
  \setlength{\partopsep}{0pt}%
  \setlength{\parsep}{0pt}%
  \setlength{\itemsep}{0pt}%
}}
\begin{Verbatim}[breaklines=true,formatcom=\color{ada_blue}]
# Create a data frame for the true line over the extended range
true_line <- data.frame(x = x_extended, y = beta_0 + beta_1 * x_extended)

# Calculate the standard error of the true regression line
mean_x <- mean(x_values) # Mean of the x values
S_xx <- sum((x_values - mean_x)^2) # Sum of squares of x deviations
se_fit_true <- sigma * sqrt(1 / n + (x_extended - mean_x)^2 / S_xx) # Standard error formula for regression

# Calculate the theoretical confidence intervals based on the true model
t_value <- qt(1 - alpha / 2, df = n - 2) # Critical t-value for 95% CI
ci_upper_true <- (beta_0 + beta_1 * x_extended) + t_value * se_fit_true # Upper bound
ci_lower_true <- (beta_0 + beta_1 * x_extended) - t_value * se_fit_true # Lower bound

# Create a data frame for the confidence interval bounds based on the true regression line
ci_bounds_true <- data.frame(x = x_extended, ci_lower = ci_lower_true, ci_upper = ci_upper_true)
\end{Verbatim}
\endgroup
\par\addvspace{\topsep}

\noindent Next, we sample $n$ points from a normal distribution, to generate $y$ values for each $x$ value and fit the corresponding regression line. We repeat this process \ttblue{rep} times.
